% !TEX program = pdflatex
% =====================================================================
% tarefas-dados-en.tex -- standalone, landscape, en-US: data table,
% method table, references
%
% Crosses the task as the energy literature defines it with the open
% Brazilian source that serves it as input or as ground truth.
%
% SCOPE. Mapping: task from the literature, open national source that
% serves it, and what that source carries. State of this repository --
% what was downloaded, what was run, which experiment closed -- does not
% enter. Survey of Aug 30 -- Sep 7, 2026.
%
% TWO DELIBERATE CHOICES FOR en-US:
%
%   `babel` loads `american`, and `\sisetup` overrides the style sheet's
%   decimal comma and thin-space grouping with the en-US period and
%   comma. The override is local to this file; `estilo-reporte.sty` keeps
%   serving the Portuguese documents unchanged.
%
%   Constrained-off, endpoint, host and bucket are set roman. An italic
%   marks a word as foreign to the surrounding text; in English these
%   are not.
%
% Dataset slugs, agency names and the CCEE error string stay in
% Portuguese: they are identifiers, and translating them breaks the
% lookup they exist to support.
%
% DATA PROPERTIES that, ignored, lead to incorrect simplification:
%
%   `id_subsistema` carries 'N ' and 'S ' with a trailing space in the
%   EAR file. Without a strip, the pivot creates a separate column and
%   dropping missing values removes Northern and Southern hydrology.
%
%   The half-hourly CMO is dispatch-model output, not ex-post
%   measurement, and it is published BEFORE the day it refers to -- the
%   file downloaded at 10:08 on Aug 27, 2026 already carried the 48
%   half-hours of that same day. This is potential leakage for any
%   forecasting experiment over it, and separates task 11 from the rest
%   in the source column.
%
%   The hourly grid of the MCTI operating margin is rectangular: 31 days
%   in every month, and the nonexistent day arrives as zero, not empty.
%   A calendar filter is mandatory.
%
%   `recurso` is the CKAN unit: the same period enters once per format.
%   27 years of national interchange are 60 resources.
%
% NUMBERING. Task 5 was removed on Sep 7, 2026 and the 5 stays vacant,
% so a task cited by number elsewhere keeps its number.
%
% REFERENCES. The 22 entries with a DOI or arXiv id resolve in Crossref
% and arXiv, checked Sep 7, 2026; `ieapvps2023` and `eprs2025` have no
% DOI. No entry supports a number in the data column, which comes
% from the source; each anchors the task as distinct in the literature,
% except `rolnick2022` and `liao2022`, which anchor the column selection
% of the second table.
%
% Change history: in git.
%
% Compiles on its own, and uses the style sheet of this directory for
% visual identity only:
%   latexmk -pdf tarefas-dados-en.tex
% =====================================================================

% The style sheet presupposes `book`: it uses `\chaptertitlename` and the
% chapter commands of `tocloft`. Under `article` it does not load.
\documentclass[10pt,a4paper,oneside]{book}

\usepackage[american]{babel}
\usepackage[a4paper,top=1.9cm,bottom=1.9cm,left=2.0cm,right=2.2cm]{geometry}
\usepackage{estilo-reporte}

% The style sheet sets a decimal comma and thin-space grouping, which is
% what the Portuguese documents need. Overridden here, after the sheet
% loads, so `\num` prints 19,088,880 and not 19 088 880.
\sisetup{output-decimal-marker={.},group-separator={,}}

% Core and auxiliary are distinguished by FILL, not by size. The two
% earlier markers were solid squares of 0.46 and 0.26 em: `\tabfont` runs
% at 7.4 pt, so the smaller one came out 1.9 pt on a side and telling the
% two apart became a matter of looking closely, cell by cell. Solid
% against outlined reads at a distance, with the same 0.55 em footprint
% in both -- what changes is the interior, not the area, and the columns
% line up. Same hue in both: what separates role here is shape, and the
% style sheet already uses lightness for tier.
\newcommand{\nuc}{\textcolor{destaque}{\raisebox{0.02em}{\rule{0.55em}{0.55em}}}}
\newcommand{\aux}{\raisebox{0.02em}{\begingroup
  \setlength{\fboxsep}{0pt}\setlength{\fboxrule}{0.12em}%
  \textcolor{destaque}{\fbox{\rule{0pt}{0.31em}\rule{0.31em}{0pt}}}\endgroup}}
\newcommand{\x}{}
% A CKAN slug does not break on its own inside `\texttt`: the monospaced
% face comes with `\hyphenchar` switched off, and the whole name spills
% over the neighboring column. This command opens a break point after
% each hyphen, without duplicating the hyphen.
% Tier: how many of the two published priority sources name the task. The
% ramp is a single hue in three lightnesses -- the style sheet
% distinguishes role by lightness and not by hue, and tier follows the
% same rule.
\definecolor{tierI}{HTML}{30401A}
\definecolor{tierII}{HTML}{658637}
\definecolor{tierIII}{HTML}{B7C79C}
\newcommand{\tchip}[3]{\colorbox{#1}{\textcolor{#2}{\fontsize{6.4}{6.4}\selectfont\bfseries\strut\,#3\,}}}
\newcommand{\tum}{\tchip{tierI}{white}{I}}
\newcommand{\tdois}{\tchip{tierII}{white}{II}}
\newcommand{\ttres}{\tchip{tierIII}{tierI}{III}}
\newcommand{\hb}{\discretionary{-}{}{-}}

% Short prose in a single column. The `promancha` of the style sheet
% opens `multicols{2}`, and in a block of three or four lines the left
% column fills mid-sentence: the eye drops to the next block instead of
% crossing to the right. Measured at the opening of this document, with
% two three-line blocks stacked. The measure is the same column as
% `promancha` -- half the text block minus half the 1.2 cm `columnsep`
% -- so the two block types align on the same line width.
% `\hsize`, and not `minipage`: the style sheet note does it this way
% because a box zeroes `\prevdepth` and the next block prints on top.
\newenvironment{promanchaum}
  {\par\begingroup
     \hsize=\dimexpr0.5\hsize-0.6cm\relax
     \linewidth=\hsize
     \parindent=0pt}
  {\par\endgroup\medskip}

% The page ends where the text ends.
\raggedbottom
\setlength{\parskip}{0.4em}
\setlength{\parindent}{0pt}

% `empty`, and not `plain`. The `plain` footer is
% centered at the foot of the UPRIGHT sheet; `pdflscape` rotates the page
% after it is composed, and the number ends up lying in the middle of the
% side margin, far from any corner. Measured: with `plain`, the "3" of
% the third page comes out turned 90 degrees, 4 cm from the left edge.
% A landscape document of this length does not need page numbers.
\pagestyle{empty}

\autoria{João Leonardi da Silva Melo}{joao\_leonardi.melo@somosicev.com}

% Closing mark, at the foot of the last page. `\fecholargura` is 45 mm,
% against the 30 mm of `\marcalargura` at the opening: the header mark
% shares its line with the title and sits near the height of one title
% line, while this one has an empty half-page to itself and reads as a
% closing device rather than as a second opening. It takes the path from
% `\re@marca`, and not the literal file, so a `\marca` in the preamble
% moves both marks.
%
% One `\vfill`, and no box after the mark. Centering it in the leftover
% space takes a second `\vfill` plus a trailing `\null`, since glue at the
% end of a page is discarded unless a box follows it -- and at 55 mm that
% `\null` no longer fitted, so it carried a blank sixth page with it. With
% a single `\vfill` the mark is pushed to the bottom margin and the count
% stays at five pages. `\raggedbottom` adds its own fill at the foot, an
% order of magnitude smaller, so it does not compete.
%
% `\vfill` alone parks the mark exactly on the bottom margin, and nothing
% below it is left to give. To sit lower, the mark has to enter the margin,
% and it does so by declaring less depth than it draws: `\raisebox` with
% the third argument at 0 pt lowers the artwork by `\fechodescida` while
% the box still measures as if it had not moved. The page builder therefore
% sees the same height as before -- no reflow, no extra page. At 8 mm of a
% 19 mm bottom margin, 11 mm of clearance remain to the paper edge, so the
% mark stays clear of the trim on print.
\newlength{\fecholargura}\setlength{\fecholargura}{45mm}
\newlength{\fechodescida}\setlength{\fechodescida}{8mm}
\makeatletter
\newcommand{\fechomarca}{%
  \par\vfill
  {\centering
   \raisebox{-\fechodescida}[\height][0pt]{%
     \includegraphics[width=\fecholargura]{\re@marca}}\par}}
\makeatother

\makeatletter
\renewenvironment{thebibliography}[1]
  {\list{\@biblabel{\@arabic\c@enumiv}}%
     {\settowidth\labelwidth{\@biblabel{#1}}%
      \leftmargin\labelwidth \advance\leftmargin\labelsep
      \usecounter{enumiv}\let\p@enumiv\@empty
      \renewcommand\theenumiv{\@arabic\c@enumiv}}%
   \sloppy\clubpenalty4000 \@clubpenalty\clubpenalty
   \widowpenalty4000 \sfcode`\.\@m}
  {\endlist}
\makeatother

\begin{document}
\begin{landscape}

\aberturaunica{Energy tasks and Brazilian open data}
              {Survey of August 30 to September 7, 2026}

\tabfont
\begin{longtable}{@{}
    >{\centering\arraybackslash}p{0.8cm}
    >{\rag\arraybackslash}p{0.55cm}
    >{\rag\arraybackslash}p{3.6cm}
    >{\rag\arraybackslash}p{5.0cm}
    >{\rag\arraybackslash}p{6.5cm}
    >{\rag\arraybackslash}p{6.5cm}@{}}
\toprule
\cab{Tier} & \cab{\#} & \cab{Task} & \cab{Stated objective} &
\cab{Open national source} & \cab{What the source provides} \\
\midrule
\endfirsthead
\toprule
\cab{Tier} & \cab{\#} & \cab{Task} & \cab{Stated objective} &
\cab{Open national source} & \cab{What the source provides} \\
\midrule
\endhead
\bottomrule
\endfoot
\tum & 2 & Generation and load forecasting &
Forecast solar generation, wind generation and load. Recent work also uses the task as a test bed for time series foundation models~\cite{sweeney2020,aksu2024} &
ONS \texttt{programacao\_x\_previsao}; \texttt{carga-energia-programada}; \texttt{curva-carga}; \texttt{demanda\_maxima\_di} &
Forecast and scheduled output per plant, in 48 daily steps, for 597 plants. The plant is identified by \texttt{cod\_usinapdp}, a third system, with no intersection with \texttt{id\_ons} or \texttt{usina\_conjunto}; matching by name, with case and whitespace normalized, resolves 26 of the 597 \\
\addlinespace[.5ex]

\tum & 10 & Asset failure and unavailability &
Estimate failure rate, availability and maintenance need of the generating unit~\cite{stetco2019} &
ONS \texttt{taxa\_teif\_teip}; \texttt{disponibilidade\_usina}; \texttt{ind\_disponibilidade\_*}; \texttt{ind-disponibilidade-geracao-sin} &
The TEIFa and TEIP rates cover 265 plants and 251 CEG on a monthly basis, from Jul 2021 to Jul 2026, with version and calculation date. Hourly availability carries installed capacity and operational availability for 158 plants identified by \texttt{id\_ons} and \texttt{ceg} \\
\addlinespace[.5ex]

\tum & 12 & Interchange, flow and network constraint &
Estimate flow between subsystems, transfer limit and energy curtailed by network constraint~\cite{bird2016,liao2022} &
ONS \texttt{intercambio-nacional}, \texttt{intercambio-internacional}, \texttt{intercambio\_modalidade}, \texttt{programacao\_fluxo\_controlado}, \texttt{restricao\_coff\_eolica\_detail} and \texttt{\_usi}, \texttt{restricao\_coff\_fotovoltaica} &
National, international and by-modality interchange cover 27 years each, from 2000 to 2026, in one resource per year and format. Controlled flow covers 700 days since 2024, and wind constrained-off, 60 months since Oct 2021, with verified generation against scheduled, plant by plant \\
\addlinespace[.5ex]

\tum & 13 & Power quality, frequency and ancillary services &
Measure frequency and voltage deviation, and the reserve that keeps them in band in a system of decreasing inertia~\cite{milano2018} &
ONS \texttt{ind\_qualid\_dfd\_evento}, \texttt{ind\_qualid\_dfd\_meano}, \texttt{ind\_qualid\_dfp\_regime}, \texttt{ind\_confiarb\_dreq\_freq}, \texttt{uge\_opera\_csi}, \texttt{usi\_servicoancilar\_sandbox}, \texttt{oferta\_respostademanda} &
The four quality and frequency indicators come in a single cumulative file, with no year in the resource name. The synchronous condenser is published since 2021, demand response since 2023, and the reactive ancillary service only in 2026 \\
\addlinespace[.5ex]

\tdois & 1 & Resource assessment &
Measure available irradiance and wind, and validate a satellite product against surface measurement~\cite{sengupta2024,ieapvps2023} &
ONS \texttt{restricao\_coff\_fotovoltaica\_detail} (plane-of-array irradiance); Rede SONDA, INPE; Atlas Brasileiro de Energia Solar, LABREN; NASA POWER &
Plane-of-array irradiance has \num{19088880} records from 560 plants in 10 states, every 30 min, from Apr 2024 to Sep 1, 2026, of which 16.9\% arrive flagged as suspect. The Atlas carries \num{72272} records of monthly and annual means \\
\addlinespace[.5ex]

\tdois & 3 & Asset inventory and detection &
Locate and date plants and distributed installations from satellite imagery~\cite{kruitwagen2021,yu2018,li2025} &
ANEEL SIGA; ANEEL MMGD; ONS \texttt{linha-transmissao} and \texttt{subestacao}; EPE; Sentinel-2 &
SIGA carries \num{25133} projects with coordinate, capacity and commissioning date. MMGD carries \num{4678147} connections, whose coordinate is that of the substation and not of the installation. The network has \num{2208} lines and \num{1677} substations \\
\addlinespace[.5ex]

\tdois & 4 & Satellite monitoring of operation &
Estimate generation and emissions of a plant without telemetry, calibrating a plume or thermal signature against reported generation~\cite{nassar2017} &
ONS \texttt{geracao-usina-2}; \texttt{geracao-termica-despacho-2}; \texttt{cvu-usitermica}; \texttt{balanco\_dessem\_detalhe} &
Thermal dispatch comes out by reason, on a monthly basis, and unit variable cost per plant, on an annual basis since 2020. The DESSEM balance carries 192 rows per day: four subsystems in 48 steps \\
\addlinespace[.5ex]

\tdois & 9 & System adequacy and reliability &
Estimate the risk of load not served and the severity of interruption, and measure unserved energy~\cite{ghanbarzadeh2025} &
ONS \texttt{ind\_confiarb\_energianaosuprida}; \texttt{interrupcao\_carga}; \texttt{energia-vertida-turbinavel}; \texttt{ind\_confiarb\_*} (robustness, severity, CIPER); \texttt{ear-diario-por-subsistema} &
Unserved energy has \num{37592} records from Jan 2007 to Jul 2026, in 32 aggregation attributes. Load interruption has \num{8325} events and \num{6821} disturbances from Jan 2007 to Sep 2026, with interrupted load, mean duration and unserved energy \\
\addlinespace[.5ex]

\tdois & 11 & Price and marginal operating cost &
Forecast the half-hourly marginal operating cost and the settlement price, with the metric used in the price forecasting literature~\cite{lago2021} &
ONS \texttt{cmo-semanal}, \texttt{cmo\hb semi\hb horario}, \texttt{ofertapreco-importacao}, \texttt{cvu-usitermica}; CCEE, Dados Abertos (PLD) &
The weekly CMO covers 2005 to 2026; the half-hourly, 2020 to 2026, for four subsystems in 48 half-hours per day; and the import price offer, 2019 to 2026. The half-hourly series is dispatch-model output, not ex-post measurement: it is published before the day it refers to \\
\addlinespace[.5ex]

\tdois & 15 & Methane leakage and fugitive emissions &
Detect and quantify a methane point plume from satellite, with detection limit in kg/h as the metric~\cite{mohammadimanesh2025} &
Sentinel-5P/TROPOMI, Sentinel-2 and EMIT, open; ANP, oil and gas production per well &
ANP publishes monthly production per well, in 54 files. It is the only national source that locates the installation as a plume candidate \\
\addlinespace[.5ex]

\tdois & 16 & Right-of-way vegetation and line risk &
Measure vegetation encroachment on the transmission corridor and the associated ignition risk~\cite{gazzea2022} &
ONS \texttt{linha-transmissao}; MapBiomas; INPE, Programa Queimadas &
INPE publishes active fire detections in 23 annual files, from 2003 to 2025, for the reference satellite, besides the 10 min, daily and monthly series. MapBiomas publishes annual land cover classification for the whole country, in raster. Line geometry comes from ONS, in \num{2208} alignments \\
\addlinespace[.5ex]

\ttres & 6 & Plant performance and degradation &
Measure performance ratio, soiling and unavailability over the life of the asset~\cite{jordan2013,jordan2022} &
ONS \texttt{restricao\_coff\_fotovoltaica\_detail} (POA, verified and estimated generation); \texttt{disponibilidade\_usina}; \texttt{taxa\_teif\_teip} &
POA, verified generation and estimated generation come out together every 30 min, from Apr 2024 to Sep 1, 2026, which gives the numerator and the denominator of the performance ratio in the same record. Availability and the unavailability rates cover Jul 2021 to Jul 2026 \\
\addlinespace[.5ex]

\ttres & 7 & Siting and land use conflict &
Select a site and measure competition with agriculture, including agrivoltaic arrangements~\cite{barrongafford2019,dunnett2020} &
MapBiomas; IBGE; ANEEL SIGA; ONS basin outlines &
MapBiomas gives the land use class in a national annual raster; SIGA, the coordinate and capacity of \num{25133} projects; and ONS, the basin outline. Crossing site against current use is immediate, and against past use as well, through the annual series \\
\addlinespace[.5ex]

\ttres & 8 & Energy access and demand &
Estimate electrification and demand where there is no metering, from night lights and administrative records~\cite{falchetta2019} &
ANEEL MMGD by municipality; ONS \texttt{curva-carga}; VIIRS &
MMGD carries the IBGE municipality code in the \num{4678147} connections, which gives a count per municipality without going through coordinates \\
\addlinespace[.5ex]

\ttres & 14 & Marginal emissions and carbon intensity &
Attribute emissions to the marginal MWh, that is, to the generator that responds to a change in load, rather than to the system average~\cite{rolnick2022,novan2024} &
ONS \texttt{geracao-usina-2}, \texttt{cvu-usitermica}, \texttt{cmo\hb semi\hb horario}, \texttt{carga-energia}; MCTI/SIRENE, Operating Margin Emission Factor &
MCTI publishes the Operating Margin by the Dispatch Data Analysis Method on an hourly basis, with base years from 2006 to 2026. The grid is rectangular: 31 days in every month, and the nonexistent day arrives as zero, not empty. The input is entirely within ONS, and the target outside it \\
\end{longtable}

\notas{Two priority sources name tasks: the taxonomy
of~\cite{rolnick2022}, which marks a task as \emph{High Leverage}, and the
European Parliament survey on artificial intelligence in the energy
sector~\cite{eprs2025}. \tum\ task named by both sources;
\tdois\ task named by one of the two; \ttres\ task named by neither.
The table runs from tier I to tier III.
Resource is the CKAN unit: the same period enters once per format, and 27
years of national interchange amount to 60 resources.}

\clearpage
{\sffamily\bfseries\large Method classes, and where each one applies\par}
\vspace{0.4em}

% The key comes BEFORE the matrix, and as a LIST. It has already been
% below the table, and later above it but inside `\notas`, in a running
% paragraph together with the citation of the column selection and the
% panel vs. learning distinction: to know what the solid square meant one
% had to read three lines of 6.8 pt prose and hunt for the word. One line
% per mark, with the mark in the first column, reads at a glance and
% competes with no note. Notes that are not the key stay below the table,
% which is where a note belongs.
\begingroup\sffamily\fontsize{7.4}{11}\selectfont\color{tinta2}
\begin{tabular}{@{}c@{\hspace{0.7em}}l@{\hspace{1.6em}}l@{}}
\nuc & \textbf{core}      & the estimand is defined in the terms of that method: without it, the task does not exist \\
\aux & \textbf{auxiliary} & the method enters as input, calibration or constraint \\
     & \textbf{empty}     & the method does not apply to the task as posed \\
\end{tabular}
\endgroup

\vspace{0.9em}
\tabfont
\begin{tabular}{@{}
    >{\centering\arraybackslash}p{0.8cm}
    >{\rag\arraybackslash}p{0.55cm}
    >{\rag\arraybackslash}p{3.3cm}
    *{8}{>{\centering\arraybackslash}p{2.05cm}}@{}}
\toprule
& & & \multicolumn{4}{c}{\cab{Problem structure}}
  & \multicolumn{1}{c}{\cab{Data}}
  & \multicolumn{3}{c}{\cab{Solution class}} \\
\cmidrule(lr){4-7}\cmidrule(lr){8-8}\cmidrule(lr){9-11}
\cab{Tier} & & & \cab{Time series} & \cab{Tabular panel} & \cab{Spatial} & \cab{Graph}
  & \cab{Computer vision}
  & \cab{Deterministic physics} & \cab{Machine learning} & \cab{Optimization} \\
\midrule
\tum & 2 & Forecasting                & \nuc & \nuc & \aux & \x   & \aux & \aux & \nuc & \x   \\
\tum & 10 & Failure and unavailability & \nuc & \nuc & \x   & \aux & \aux & \aux & \nuc & \aux \\
\tum & 12 & Interchange and flow       & \nuc & \aux & \aux & \nuc & \x   & \nuc & \aux & \nuc \\
\tum & 13 & Quality and frequency      & \nuc & \aux & \x   & \aux & \x   & \nuc & \aux & \aux \\
\tdois & 1 & Resource assessment        & \aux & \aux & \nuc & \x   & \x   & \nuc & \aux & \x   \\
\tdois & 3 & Inventory and detection    & \nuc & \x   & \nuc & \aux & \nuc & \x   & \nuc & \x   \\
\tdois & 4 & Satellite monitoring       & \aux & \nuc & \aux & \x   & \nuc & \aux & \nuc & \aux \\
\tdois & 9 & Adequacy and reliability   & \nuc & \aux & \aux & \nuc & \x   & \nuc & \aux & \nuc \\
\tdois & 11 & Price and marginal cost    & \nuc & \aux & \aux & \aux & \x   & \aux & \nuc & \nuc \\
\tdois & 15 & Methane and leakage        & \aux & \x   & \nuc & \x   & \nuc & \nuc & \nuc & \x   \\
\tdois & 16 & Right-of-way and risk      & \aux & \x   & \nuc & \aux & \nuc & \x   & \nuc & \aux \\
\ttres & 6 & Performance, degradation  & \nuc & \nuc & \aux & \x   & \x   & \nuc & \aux & \x   \\
\ttres & 7 & Siting and land use        & \aux & \x   & \nuc & \aux & \nuc & \aux & \aux & \nuc \\
\ttres & 8 & Access and demand          & \aux & \nuc & \nuc & \x   & \nuc & \x   & \nuc & \x   \\
\ttres & 14 & Marginal emissions         & \nuc & \nuc & \aux & \aux & \x   & \aux & \nuc & \aux \\
\bottomrule
\end{tabular}

\notas{The column selection follows~\cite{rolnick2022}, with the graph in a
column of its own because it constitutes there a method class with its own
formulation and metrics~\cite{liao2022}, and not a particular case of a panel.
\emph{Tabular panel} is the structure of the data, plant $\times$ time, and
\emph{machine learning} is the technique that solves it: a panel can be fitted
by a declared regression, and machine learning also operates over imagery and
over series.}

% `multicols` zeroes `\prevdepth`, so the glue between the table and the
% block disappears and the first line of prose prints over the
% `\bottomrule`. Measured here, with the note moved above the table:
% without this glue, the row for task 14 and the opening sentence occupy
% the same height. The space goes OUTSIDE `promancha`; a `\vspace` inside
% it falls into the column, and not before the block.
\par\vspace{1.1em}

\begin{promancha}
Nine tasks have a time series core. In five of them, the estimand is already
published as a series: performance (6), with POA and generation at 30 min since
Apr 2024; adequacy (9) and asset failure (10), whose entire history fits in one
file, with unserved energy and load interruption since Jan 2007 and TEIFa and
TEIP since Jul 2021; price (11), with the half-hourly CMO since 2020; and
marginal emissions (14), with the hourly operating margin since 2006.

In the other four the estimand requires assembly, and the obstacle differs. In
forecasting (2) it is linkage, through the identifier proper to the scheduling
file. In asset detection (3) it is imagery: the national source gives the
inventory, not the series that would date it. In network flow (12) and in
frequency quality (13) it is format, with interchange in one resource per year
and the frequency indicators in a single cumulative file, with no year in the
name.

None of the fifteen tasks lacks published national data. Task 14 is the only case
with the input in ONS and the estimand outside it, and the estimand is published
as well: the marginal emission factor comes from MCTI on an hourly basis since
2006.
A single input serves three tasks. Plane-of-array irradiance from
560 plants, every half hour, is an input to resource assessment, ground truth
for forecasting and an input to the performance ratio. There is no public
equivalent in the country: Rede SONDA measures in the horizontal plane, in a
network of dozens of stations, and the LABREN Atlas is a modeled product in
monthly means.
\end{promancha}

\clearpage
{\sffamily\bfseries\large Open national sources outside ONS\par}
\vspace{0.4em}

\begin{promanchaum}
CCEE is independent of the operator, a condition of task 4, which compares a
satellite estimate against the reported value: ONS on both sides does not
constitute an independent comparison. The remaining sources cover what ONS does
not publish: interruption at the distribution level, surface meteorological
measurement, active fire detections and land cover.
\end{promanchaum}

\vspace{0.5em}
\tabfont
\begin{tabular}{@{}
    >{\rag\arraybackslash}p{4.4cm}
    >{\rag\arraybackslash}p{5.1cm}
    >{\rag\arraybackslash}p{7.4cm}
    >{\rag\arraybackslash}p{1.3cm}
    >{\rag\arraybackslash}p{5.0cm}@{}}
\toprule
\cab{Agency and dataset} & \cab{What it provides} & \cab{Coverage and volume} &
\cab{Tasks} & \cab{Access route} \\
\midrule

ANEEL, Dados Abertos: \texttt{interrupcoes\hb de\hb energia\hb eletrica\hb nas\hb redes\hb de\hb distribuicao} &
Every interruption in the country's distribution network, event by event, outside the areas of permit-holding utilities &
One file per year, from 2017 to 2026, in ZIP and Parquet, from 78 MB (2017) to 272 MB (2025) &
9 &
Open CKAN API, no registration \\
\addlinespace[.5ex]

ANEEL, Dados Abertos: \texttt{indicadores\hb coletivos\hb de\hb continuidade\hb dec\hb e\hb fec} &
DEC and FEC per set of consumer units, with the compensation paid for violation, the regulatory limit and the physical-electrical attributes of the set &
Three ranges, 2000--2009, 2010--2019 and 2020--2029, besides compensation, limit, set attributes and domain table &
9 &
Open CKAN API, no registration \\
\addlinespace[.5ex]

INMET, automatic stations and BDMEP &
Hourly surface measurement: temperature, humidity, wind, precipitation and global radiation &
674 automatic stations, 518 operating and 156 out of service, distributed over 25 states; RR and SE are absent. The oldest has operated since May 6, 2000 &
1, 2 &
Catalog through the API, no registration. The hourly endpoint of \texttt{apitempo} returns 204 No Content and does not serve the series; global radiation comes through the annual ZIP of \texttt{portal.inmet.gov.br} \\
\addlinespace[.5ex]

CCEE, Dados Abertos &
Generation and consumption from market settlement, per plant and per source, previously released in the InfoMercado spreadsheets &
No measured coverage: the API responds HTTP 403, ``Acesso bloqueado'' &
2, 4, 11, 14 &
The block applies to all three routes: CKAN direct, CKAN with a \texttt{Referer} from the portal itself, and the host \texttt{www.ccee.org.br}. The portal responds in the browser, and the programmatic route does not \\
\addlinespace[.5ex]

INPE, Programa Queimadas: active fire &
Satellite fire detections, event by event with date, time and coordinate &
Annual ZIP file from 2003 to 2025 in four extents: \texttt{Brasil\hb sat\hb ref}, \texttt{Brasil\hb todos\hb sats}, \texttt{AMS\hb sat\hb ref} and \texttt{EstadosBr\hb sat\hb ref} &
16 &
Open file server, no registration. The path \texttt{anual/Brasil/} returns 404; the four extents above are the ones that exist \\
\addlinespace[.5ex]

MapBiomas, land cover and use collection &
Annual land cover and use classification for the whole country, in raster &
Collection 11 at 30 m occupies 763 MB per year, and collection 4 at 10 m, 4.7 GB, in BigTIFF &
7, 16 &
GeoTIFF direct from the public bucket, no registration and without going through Earth Engine \\

\bottomrule
\end{tabular}

\notas{File size is the one declared by the resource, except for INPE and
MapBiomas, measured in the HTTP header.}

\begin{promancha}
\vspace{0.8em}
The target of task 14 is at MCTI. Task 9 has two scales of interruption: the
transmission scale, from ONS, and the distribution scale, event by event since
2017 and as a set indicator since 2000. INMET gives task 1 a further 518 operating
surface stations, against the dozens of Rede SONDA; they measure in the horizontal plane
and not in the plane of the array, and therefore complement the ONS POA without
replacing it. Task 16 goes from the isolated line geometry to the crossing with
the corridor, through INPE fire detections and MapBiomas land cover, both without
registration; MapBiomas also serves task 7.
\end{promancha}

\begin{objecao}
The data cell gives the coverage the source declares --- period, granularity,
entity count ---, and this establishes neither a continuous series nor a filled
column. The column name can also mislead: in
\texttt{val\_irradianciaverificado}, it diverges from the quantity the field
carries.
\end{objecao}

% No `\clearpage` here, unlike the two sections above. Those open with a
% wide `tabular`, which does not break across pages: starting mid-sheet
% would push the whole table to the next one and the gap would return.
% The bibliography breaks item by item, so it flows and fills the page of
% the objection instead of leaving it with four lines.
\vspace{1.4em}
{\sffamily\bfseries\large References\par}
\vspace{0.3em}
{\sffamily\fontsize{6.8}{8.6}\selectfont\color{tinta2}
\begin{thebibliography}{99}
\setlength{\itemsep}{0.15em}

\bibitem{sengupta2024} Sengupta, M., Habte, A., Wilbert, S., Gueymard, C.,
  Remund, J., Lorenz, E., van Sark, W. and Jensen, A.R. \emph{Best practices
  handbook for the collection and use of solar resource data for solar energy
  applications: fourth edition}. NREL/TP-5D00-88300, National Renewable Energy
  Laboratory, 2024. doi: 10.2172/2448063.

\bibitem{ieapvps2023} IEA-PVPS Task 16. \emph{Worldwide benchmark of modelled
  solar irradiance data}. Report T16-05:2023, 2023. Compares ten datasets from
  nine providers against 129 surface stations.

\bibitem{sweeney2020} Sweeney, C., Bessa, R.J., Browell, J. and Pinson, P. The
  future of forecasting for renewable energy. \emph{WIREs Energy and
  Environment}, 9(2), e365, 2020. doi: 10.1002/wene.365.

\bibitem{aksu2024} Aksu, T., Woo, G., Liu, J., Liu, X., Liu, C., Savarese, S.,
  Xiong, C. and Sahoo, D. \emph{GIFT-Eval: a benchmark for general time series
  forecasting model evaluation}. arXiv:2410.10393, 2024. Includes the energy
  domain among the seven covered.

\bibitem{kruitwagen2021} Kruitwagen, L., Story, K.T., Friedrich, J., Byers, L.,
  Skillman, S. and Hepburn, C. A global inventory of photovoltaic solar energy
  generating units. \emph{Nature}, 598, 604--610, 2021.

\bibitem{yu2018} Yu, J., Wang, Z., Majumdar, A. and Rajagopal, R. DeepSolar: a
  machine learning framework to efficiently construct a solar deployment
  database in the United States. \emph{Joule}, 2(12), 2605--2617, 2018.

\bibitem{li2025} Li, A., Liu, L., Li, S., Cui, X., Chen, X. and Cao, X.
  Global photovoltaic solar panel dataset from 2019 to 2022.
  \emph{Scientific Data}, 12, 637, 2025.
  doi: 10.1038/s41597-025-04985-y.

\bibitem{nassar2017} Nassar, R., Hill, T.G., McLinden, C.A., Wunch, D., Jones,
  D.B.A. and Crisp, D. Quantifying CO$_2$ emissions from individual power plants
  from space. \emph{Geophysical Research Letters}, 44(19), 2017.
  doi: 10.1002/2017GL074702.

\bibitem{jordan2013} Jordan, D.C. and Kurtz, S.R. Photovoltaic degradation
  rates --- an analytical review. \emph{Progress in Photovoltaics: Research and
  Applications}, 21(1), 12--29, 2013 (published online in 2011).
  doi: 10.1002/pip.1182.

\bibitem{jordan2022} Jordan, D.C., Anderson, K., Perry, K., Muller, M.,
  Deceglie, M., White, R. and Deline, C. Photovoltaic fleet degradation insights.
  \emph{Progress in Photovoltaics: Research and Applications}, 30(10),
  1166--1175, 2022. doi: 10.1002/pip.3566.

\bibitem{barrongafford2019} Barron-Gafford, G.A., Pavao-Zuckerman, M.A., Minor,
  R.L. et al. Agrivoltaics provide mutual benefits across the food--energy--water
  nexus in drylands. \emph{Nature Sustainability}, 2, 848--855, 2019.
  doi: 10.1038/s41893-019-0364-5.

\bibitem{dunnett2020} Dunnett, S., Sorichetta, A., Taylor, G. and Eigenbrod, F.
  Harmonised global datasets of wind and solar farm locations and power.
  \emph{Scientific Data}, 7, 130, 2020. doi: 10.1038/s41597-020-0469-8.

\bibitem{ghanbarzadeh2025} Ghanbarzadeh, T., Habibi, D. and Aziz, A. Addressing
  reliability challenges in generation capacity planning under high penetration
  of renewable energy resources and storage solutions: a review. \emph{Renewable
  and Sustainable Energy Reviews}, 212, 115461, 2025.
  doi: 10.1016/j.rser.2025.115461.

\bibitem{stetco2019} Stetco, A., Dinmohammadi, F., Zhao, X., Robu, V., Flynn,
  D., Barnes, M., Keane, J. and Nenadi\'c, G. Machine learning methods for wind
  turbine condition monitoring: a review. \emph{Renewable Energy}, 133,
  620--635, 2019. doi: 10.1016/j.renene.2018.10.047.

\bibitem{falchetta2019} Falchetta, G., Pachauri, S., Parkinson, S. and Byers, E.
  A high-resolution gridded dataset to assess electrification in sub-Saharan
  Africa. \emph{Scientific Data}, 6, 110, 2019.
  doi: 10.1038/s41597-019-0122-6.

\bibitem{lago2021} Lago, J., Marcjasz, G., De Schutter, B. and Weron, R.
  Forecasting day-ahead electricity prices: a review of state-of-the-art
  algorithms, best practices and an open-access benchmark. \emph{Applied
  Energy}, 293, 116983, 2021. doi: 10.1016/j.apenergy.2021.116983.

\bibitem{bird2016} Bird, L., Lew, D., Milligan, M. et al. Wind and solar energy
  curtailment: a review of international experience. \emph{Renewable and
  Sustainable Energy Reviews}, 65, 577--586, 2016.

\bibitem{milano2018} Milano, F., Dörfler, F., Hug, G., Hill, D.J. and Verbič, G.
  Foundations and challenges of low-inertia systems. \emph{2018 Power Systems
  Computation Conference (PSCC)}, 1--25, 2018.
  doi: 10.23919/PSCC.2018.8450880.

\bibitem{mohammadimanesh2025} Mohammadimanesh, F., Mahdianpari, M., Radman, A.,
  Varon, D., Hemati, M. and Marjani, M. Advancements in satellite-based methane
  point source monitoring: a systematic review. \emph{ISPRS Journal of
  Photogrammetry and Remote Sensing}, 224, 94--112, 2025.
  doi: 10.1016/j.isprsjprs.2025.03.020.

\bibitem{gazzea2022} Gazzea, M., Pacevicius, M., Dammann, D.O., Sapronova, A.,
  Lunde, T.M. and Arghandeh, R. Automated power lines vegetation monitoring using
  high-resolution satellite imagery. \emph{IEEE Transactions on Power Delivery},
  37(1), 308--316, 2022. doi: 10.1109/TPWRD.2021.3059307.

\bibitem{novan2024} Novan, K. and Wang, Y. Estimates of the marginal curtailment
  rates for solar and wind generation. \emph{Journal of Environmental Economics
  and Management}, 2024.

\bibitem{eprs2025} European Parliamentary Research Service (EPRS).
  \emph{Artificial intelligence and the energy sector}. Briefing
  PE 775.859, European Parliament, 2025. Names, for the electricity sector,
  consumption and production forecasting, variable renewable integration,
  congestion management, stabilization and storage, predictive maintenance,
  anticipation of disruption by extreme event, demand response and price
  forecasting.

\bibitem{rolnick2022} Rolnick, D., Donti, P.L., Kaack, L.H., Kochanski, K.,
  Lacoste, A., Sankaran, K. et al. Tackling climate change with machine
  learning. \emph{ACM Computing Surveys}, 55(2), article 42, 1--96, 2022.
  doi: 10.1145/3485128. Anchors the column selection of the second table,
  not a task.

\bibitem{liao2022} Liao, W., Bak-Jensen, B., Pillai, J.R., Wang, Y. and Wang, Y.
  A review of graph neural networks and their applications in power systems.
  \emph{Journal of Modern Power Systems and Clean Energy}, 10(2), 345--360,
  2022. doi: 10.35833/MPCE.2021.000058.

\end{thebibliography}
}

\fechomarca

\end{landscape}
\end{document}
